\section{Discusión de \textit{Trade-offs}}
\label{sec:tradeoffs}

\subsection{Quantum fijo vs. bloque proporcional: costo de coordinación}

Los modos cooperativo y quantum son mecánicamente idénticos --- comparten
el mismo bucle del worker (Algoritmo~\ref{alg:worker}) y la misma función de
ejecución; solo difiere cómo se calcula el tamaño del bloque por activación
(Sección~\ref{sec:algoritmos}). Esa única diferencia tiene, sin embargo, un
costo de coordinación medible y opuesto según el tamaño relativo de las
tareas. Con quantum fijo $Q$, una tarea pequeña frente a $Q$ puede terminar
en una sola activación mientras una tarea grande necesita muchas; con un
bloque proporcional al tamaño de cada tarea, todas necesitan
aproximadamente el mismo número de activaciones sin importar su
\texttt{work\_units} total, a costa de que una tarea grande ceda el CPU con
más frecuencia de la que cedería bajo un quantum fijo equivalente. El Caso~6
del enunciado lo hace explícito con una sola tarea de 60 unidades: al 10\,\%
cooperativo (bloques de 6) necesita 10 activaciones; con quantum $Q=25$
(bloques 25+25+10) necesita solo 3, para exactamente el mismo trabajo total
y el mismo resultado de $\pi$. Más activaciones significan más ciclos de
mutex/condvar y más filas de log por unidad de trabajo real completada ---
el \textit{overhead} de coordinación del protocolo de la
Sección~\ref{sec:protocolo}, que no aparece en el resultado final observable
pero sí en el costo de producirlo.

\subsection{Reproducibilidad vs. realismo temporal}

El enunciado exige que la misma entrada, semilla, modo y parámetros
produzcan exactamente el mismo registro de decisiones (determinismo). Esto
se logra midiendo el trabajo en \texttt{work\_units} lógicas --- iteraciones
discretas de la serie de $\pi$ --- en vez de tiempo real de CPU, y usando
aritmética entera (Algoritmo~\ref{alg:cooperativo}) en vez de \texttt{double}
para el cálculo del bloque cooperativo, evitando que la representación de
punto flotante introduzca variación entre corridas idénticas por diferencias
de redondeo entre plataformas. El costo de esa reproducibilidad es
\textbf{realismo}: en un sistema real, el tiempo que una tarea recibe
depende de interrupciones de temporizador, latencia de memoria, otros
procesos del sistema operativo real corriendo por debajo, y variabilidad de
\textit{hardware} --- ninguno de esos factores está presente aquí. La
elección es deliberada y coherente con el objetivo del proyecto (verificar
un protocolo de sincronización y una política de asignación, no medir
rendimiento real de un sistema), pero implica que los números de
\textit{overhead} de coordinación de la sección anterior son comparables
\emph{entre sí} (mismo entorno lógico) y no directamente extrapolables al
costo real de cambio de contexto de un \textit{scheduler} de producción.

\subsection{Selección $O(n)$ por lista vs. árbol de sumas parciales $O(\log n)$}

El sorteo (Algoritmo~\ref{alg:sorteo}) localiza a la ganadora recorriendo la
lista de tareas \texttt{READY} y acumulando boletos hasta que el boleto
sorteado cae en el rango de alguna --- costo $O(n)$ por despacho, donde $n$
es el número de tareas activas. La Opción E del enunciado
(Tabla~\ref{tab:opciones}) propone sustituir esa lista por un árbol de sumas
parciales, reduciendo la localización a $O(\log n)$. Para el rango de
tareas que exige este proyecto (5 a 25), la diferencia entre $O(n)$ y
$O(\log n)$ es de a lo sumo un factor de $25/\log_2 25 \approx 5$
comparaciones por despacho --- indistinguible en la práctica frente al costo
fijo de las dos llamadas a \texttt{pthread\_mutex\_lock}/\texttt{unlock} y
la señal de condvar que rodean cada sorteo. La complejidad adicional de
mantener un árbol balanceado (actualizar sumas parciales cada vez que una
tarea cambia de estado o de boletos efectivos, en vez de solo recorrer un
arreglo) solo se justificaría con cientos o miles de tareas activas
simultáneas, un régimen muy distinto al de este proyecto; de ahí que se
prefiriera reutilizar el mecanismo cooperativo/quantum (Opción A) en vez de
optimizar una operación que no es el cuello de botella real a esta escala.

\subsection{Equidad de corto plazo vs. largo plazo}

La proporcionalidad de la lotería es, por diseño, una propiedad de largo
plazo (Sección~\ref{sec:analisis-extension}): si una corrida se deja
avanzar hasta que todas las tareas terminan, \texttt{observed\_share}
converge a $1/N$ para todas, sin importar sus boletos, simplemente porque
cada tarea termina consumiendo exactamente su trabajo total y las que
tienen más boletos solo llegan ahí más rápido. La ponderación real por
boletos únicamente es observable dentro de una \textbf{ventana truncada}
(\texttt{--max-dispatches}) sobre trabajo suficientemente grande, como en
los Casos 3 y 4 (Sección~\ref{sec:resultados}). Esto implica un
\textit{trade-off} de medición, no solo de diseño: la elección del tamaño de
la ventana observada determina qué se puede concluir sobre equidad ---una
ventana demasiado corta introduce más varianza de corto plazo (menos
sorteos acumulados para que la ley de los grandes números converja);
demasiado larga, y algunas tareas empiezan a terminar, sesgando el share de
las que quedan. La elección de ventana usada en los Casos 3/4 de este
informe ($10^4$ despachos sobre tareas de $10^7$ unidades) se validó
verificando que ninguna tarea llega a completar su trabajo dentro de la
ventana (Sección~\ref{sec:resultados}); la sensibilidad de este mismo
\textit{trade-off} es aún más pronunciada al medir el efecto de la
extensión de compensación, cuyo análisis de robustez frente al tamaño de
ventana está en elaboración (Sección~\ref{sec:resultados-extension}).
